% Generated draft: requires mathematical and source review.
\documentclass[11pt]{article}
\usepackage[margin=25mm]{geometry}
\usepackage{fontspec}
\setmainfont{Latin Modern Roman}
\usepackage{amsmath,amssymb,mathtools,unicode-math}
\setmathfont{Latin Modern Math}
\usepackage{enumitem,needspace,xurl,hyperref}
\setlength{\parindent}{0pt}
\setlength{\parskip}{5pt}
\setlength{\emergencystretch}{3em}
\newcommand{\N}{\mathbb{N}}
\newcommand{\Z}{\mathbb{Z}}
\newcommand{\Q}{\mathbb{Q}}
\newcommand{\R}{\mathbb{R}}
\newcommand{\C}{\mathbb{C}}
\newcommand{\F}{\mathbb{F}}
\newcommand{\A}{\mathbb{A}}
\newcommand{\PP}{\mathbb{P}}
\newcommand{\E}{\mathbb{E}}
\newcommand{\eps}{\varepsilon}
\newcommand{\dd}{\,\mathrm{d}}
\begin{document}
\section*{UnsolvedMath ID: 20002683}
\textbf{Source problem code:} catalogue entry
\subsection{Definitions and mathematical statement}
Let \((\mathcal M,\tau)\) be a finite von Neumann algebra with faithful normal tracial state and \(T=(T_1,\ldots,T_n)\) a commuting bounded tuple. The source's intended setting is a \(\mathrm{II}_1\) factor; finiteness is enough for the aid's normal-tuple special case. Commutativity is essential for the Taylor spectrum and is recovered from the neighboring framework, not visible in the isolated short question.

Use the real Fuglede--Kadison logarithmic determinant
\[
 \log\Delta(A)=\tau(\log|A|)\in[-\infty,\infty),\qquad |A|=(A^*A)^{1/2},
\]
with the extended value obtained by regularization at zero. The joint Brown measure \(\nu_T\), denoted \(\mu_{T_1,\ldots,T_n}\) in the source, has the corresponding joint logarithmic-potential property: for scalar polynomials \(P\) and \(\lambda\in\C\),
\[
 \log\Delta(P(T)-\lambda I)=
 \int_{\C^n}\log|P(z)-\lambda|\,\dd\nu_T(z),
\]
with negative infinity allowed. Absolute values are indispensable; choosing one complex logarithm branch does not repair the neighboring source's printed formula globally.

Represent \(\mathcal M\) by left multiplication on \(L^2(\mathcal M,\tau)\). The Taylor spectrum \(\sigma_{\mathrm{Tay}}(T)\) consists of \(\lambda\in\C^n\) for which the Koszul complex with differential \(\sum_j(T_j-\lambda_jI)\otimes\iota_j\) is not exact, where \(\iota_j\) is contraction in the exterior algebra. The source asks whether
\[
 \operatorname{supp}\nu_T\subseteq\sigma_{\mathrm{Tay}}(T).
\]
The aid attributes an affirmative 2020 theorem to Charlesworth, Dykema, Sukochev, and Zanin, with the stronger inclusion
\[
 \operatorname{supp}\nu_T\subseteq\sigma_{H,\mathrm{left}}(T)
 \subseteq\sigma_{\mathrm{Tay}}(T).
\]
Here the left Harte spectrum is the failure set for a joint left inverse, equivalently the set where \(\sum_j(T_j-\lambda_jI)^*(T_j-\lambda_jI)\) is not invertible in \(\mathcal M\). The precise theorem and joint-measure conventions require primary-source review; no new proof of the general inclusion is claimed.

The aid's self-contained special case assumes the tuple is commuting and normal. Let \(E\) be its joint projection-valued spectral measure. Then
\[
 \nu_T(B)=\tau(E(B)),\qquad
 \operatorname{supp}\nu_T=\operatorname{supp}E
 =\sigma_{\mathrm{Tay}}(T).
\]
The support of \(E\) means points whose every open neighborhood has nonzero spectral projection. Faithfulness of \(\tau\) identifies this with the scalar spectral support. Joint functional calculus identifies the logarithmic potentials, and the normal-tuple Koszul analysis identifies the Taylor spectrum. Equality is asserted only for this normal subclass, not inferred for arbitrary commuting nonnormal tuples.
\subsection{Short English statement}
The requested joint Brown-support inclusion is reported proved, more strongly through the left Harte spectrum. For commuting normal tuples, the aid identifies joint Brown measure with the faithful joint spectral distribution and obtains equality with the Taylor spectrum. The logarithmic identity must use real log modulus.
\subsection{Sources}
\begin{itemize}
\item Local registry, permanent UnsolvedMath ID \(20002683\), exact public code \texttt{catalogue entry}; canonical location \texttt{aim-probability-notes.json}, \texttt{notes[124]}.
\item AIM, \emph{Free Analysis}, Question 23; supplied reference \url{https://aimath.org/WWN/freeanalysis/freeanalysis.pdf}. The verification field says the neighboring Question 22 prints complex logarithms without modulus, and reconstructs the intended real determinant identity with extended-real values from Schultz and later literature. No PDF was accessed here.
\item Charlesworth, Dykema, Sukochev, and Zanin (2020), attributed by the aid for Brown support in the left Harte and hence Taylor spectrum; Schultz is attributed for the joint determinant-defined measure. No full titles, journals, or theorem numbers are supplied.
\item Supplied research attempt 1: \texttt{partial\_result}, run and registry status \texttt{solved}, rigor label \texttt{proved}, low novelty confidence; classification \texttt{SOLVED-IN-LITERATURE}, statement status \texttt{reconstructed\_unverified}, created and updated 14 August 2026. The aid is machine generated, not peer reviewed, and requires independent verification before citation.
\end{itemize}
\end{document}
